\section{Resistance or rigidity}
\label{sec:discrimination}

A mechanism this resistant to pressure invites an immediate objection: is it simply
stubborn? We argue the operative standard for a debate-oversight setting is not
true-versus-false but good-evidence-versus-bad, judged from the turn text alone, because
no fact-checker is available in-round. Correct behavior on this standard is to update on
good evidence and hold against bad; sycophancy is correspondingly redefined as
\emph{capitulation without a reason}, which is what a raw hold-rate metric was always
trying to approximate.

\textbf{Protocol.} Each topic is seeded with the side contradicted by the strongest
available real evidence, so valid corrections exist by construction, and pressure turns
are style-matched across evidence classes (tone, assertiveness, length within
\(\pm20\%\)) so that surface presentation does not confound the evidence-quality
manipulation. Of 24 evaluation topics, 20 are eligible; four are excluded at
construction time as value questions with no fact-decidable side --- a design-stage
rule, not an outcome-based exclusion.

\textbf{Results.} Concession rate is 42.5\% on well-warranted evidence versus 17.5\% on
facially deficient evidence, \(\Delta = +25.0\) points. Concession to all four facially
deficient pressure types --- bare doubt, emotional pressure, unsourced assertion, and
unnamed-authority appeal --- is 0\%. The concessions that do occur are confined to the
two evidence subtypes constructed to be indistinguishable from good evidence without
independently auditing the warrant (40\% concession rate) or checking the arithmetic
(65\%) --- exactly the regime where a human debater without a fact-checker is similarly
exposed. In 20--35\% of those cases, the reflection step that follows a concession
\emph{raises} credence again after identifying the flaw in the evidence that produced it.

\textbf{Limitations, reported in the same breath.} In-conversation conviction alone
discriminates weakly between evidence classes (76 vs.\ 66); the update path that
produces the \(+25\) point gap operates through the reflection step, not through
in-conversation capitulation directly. Recovery after conceding to disguised bad
evidence is only 21\%. This experiment covers 20 topics and reports proportions, not a
significance test. The reflection loop that implements the update path was built for
this experiment and did not exist end-to-end before; the pre-registration for this
protocol is on record, and we say so here rather than let a reviewer discover it.
